\documentclass[9pt,a4paper,twocolumn]{ctexart}

% --- Packages ---
\usepackage[top=1.2cm, bottom=1.2cm, left=1.2cm, right=1.2cm, columnsep=0.8cm]{geometry}
\usepackage{hyperref}
\usepackage{graphicx}
\usepackage{float}
\usepackage{longtable}
\usepackage{tabularx}
\usepackage{booktabs}
\usepackage{enumitem}
\usepackage{xcolor}
\usepackage{fancyhdr}
\usepackage{listings}
\usepackage{fontspec}
\usepackage{titlesec}
\usepackage{caption}

% --- Code block styling ---
\definecolor{codebg}{RGB}{245,245,245}
\definecolor{codeframe}{RGB}{220,220,220}
\definecolor{codekw}{RGB}{0,0,192}
\definecolor{codecomment}{RGB}{0,128,0}
\definecolor{codestring}{RGB}{163,21,21}
\definecolor{codenum}{RGB}{128,0,128}
\definecolor{codepunct}{RGB}{90,90,90}
\definecolor{badgegreen}{RGB}{34,139,34}
\definecolor{badgeblue}{RGB}{30,144,255}
\definecolor{badgeblack}{RGB}{50,50,50}

% --- Difficulty badges (rounded corners via tikz, pre-rendered in saveboxes) ---
\usepackage{tikz}
\newsavebox{\badgechubox}
\savebox{\badgechubox}{\tikz[baseline]{\node[fill=badgegreen,text=white,rounded corners=2.5pt,inner xsep=4pt,inner ysep=1.5pt,font=\sffamily\footnotesize\bfseries]{初级};}}
\newsavebox{\badgezhongbox}
\savebox{\badgezhongbox}{\tikz[baseline]{\node[fill=badgeblue,text=white,rounded corners=2.5pt,inner xsep=4pt,inner ysep=1.5pt,font=\sffamily\footnotesize\bfseries]{中级};}}
\newsavebox{\badgegaobox}
\savebox{\badgegaobox}{\tikz[baseline]{\node[fill=badgeblack,text=white,rounded corners=2.5pt,inner xsep=4pt,inner ysep=1.5pt,font=\sffamily\footnotesize\bfseries]{高级};}}
\newcommand{\lvchu}{\usebox{\badgechubox}}
\newcommand{\lvzhong}{\usebox{\badgezhongbox}}
\newcommand{\lvgao}{\usebox{\badgegaobox}}
\pdfstringdefDisableCommands{%
  \def\lvchu{[初级]}%
  \def\lvzhong{[中级]}%
  \def\lvgao{[高级]}%
}

\lstset{
  basicstyle=\ttfamily\scriptsize,
  keywordstyle=\color{codekw}\bfseries,
  commentstyle=\color{codecomment}\itshape,
  stringstyle=\color{codestring},
  backgroundcolor=\color{codebg},
  frame=single,
  rulecolor=\color{codeframe},
  breaklines=true,
  breakatwhitespace=false,
  breakindent=0pt,
  postbreak=\mbox{\textcolor{red}{$\hookrightarrow$}\space},
  numbers=left,
  numberstyle=\tiny\color{gray},
  columns=flexible,
  keepspaces=true,
  showstringspaces=false,
  tabsize=2,
  aboveskip=4pt,
  belowskip=4pt,
  extendedchars=true,
  literate={，}{{，}}1 {；}{{；}}1 {！}{{！}}1 {？}{{？}}1 {“}{{“}}1 {”}{{”}}1 {（}{{（}}1 {）}{{）}}1 {《}{{《}}1 {》}{{》}}1,
}

% --- Extra language definitions (no built-in listings support) ---
\lstdefinelanguage{json}{
  morekeywords={true,false,null},
  sensitive=true,
  morestring=[b]",
  comment=[l]{//},
  morecomment=[s]{/*}{*/},
}
\lstdefinelanguage{yaml}{
  morekeywords={true,false,null,yes,no,on,off},
  sensitive=false,
  morecomment=[l]{\#},
  morestring=[b]",
  morestring=[d]',
}
\lstdefinelanguage{http}{
  morekeywords={GET,POST,PUT,DELETE,PATCH,HEAD,OPTIONS,HTTP,Host,Accept,Authorization,Content-Type,Content-Length,User-Agent},
  sensitive=true,
  morecomment=[l]{\#},
  morestring=[b]",
}

% --- Compact spacing ---
\linespread{0.95}
\setlength{\parskip}{1pt plus 1pt minus 1pt}
\setlength{\parindent}{0pt}

% --- Table styling ---
\renewcommand{\arraystretch}{1.1}

% --- Heading styling (numbered) ---
\titleformat{\section}{\normalsize\bfseries}{\thesection\quad}{0em}{}
\titlespacing{\section}{0pt}{6pt}{2pt}
\titleformat{\subsection}{\small\bfseries}{\thesubsection\quad}{0em}{}
\titlespacing{\subsection}{0pt}{4pt}{1pt}
\titleformat{\subsubsection}{\small\bfseries}{\thesubsubsection\quad}{0em}{}
\titlespacing{\subsubsection}{0pt}{3pt}{1pt}

% --- Page style ---
\pagestyle{fancy}
\fancyhf{}
\fancyhead[L]{\small\emph{\leftmark}}
\fancyhead[R]{\small\thepage}
\renewcommand{\headrulewidth}{0.4pt}
\renewcommand{\sectionmark}[1]{}  % prevent sections from overwriting leftmark

% --- Hyperref setup ---
\hypersetup{
  colorlinks=true,
  linkcolor=blue,
  urlcolor=blue,
  citecolor=blue,
}

% --- Caption format ---
\DeclareCaptionFormat{myformat}{\small\bfseries #1#2#3}
\captionsetup{format=myformat}

\begin{document}

\title{AI Agent 模拟面试题库}
\date{}
\maketitle
\markboth{Agent 知识整理}{Agent 知识整理}

\markboth{Agent 核心概念}{Agent 核心概念}


\section{1. 请用一句话概括 AI Agent 是什么？它和普通 Chatbot 有什么区别？ \lvchu{}}


\subsection{参考答案要点}

\begin{itemize}[itemsep=1pt, topsep=2pt]
  \item \textbf{Agent 定义}：Agent 是一个能感知环境、做决策、执行动作的软件系统，LLM 负责推理决策，工具负责执行，记忆负责保存上下文和历史经验。
  \item \textbf{核心公式}：Agent = LLM + Planning + Memory + Tools。
  \item \textbf{与 Chatbot 的本质区别}：
  \item Chatbot 是"你问它答"，不会主动执行操作；
  \item Agent 会在动态环境里持续观察、判断、执行，直到任务结束；
  \item Agent 能调用外部工具（API、文件系统、代码执行），不只是生成文字。
  \item \textbf{设计理念差异}：Chatbot 解决"会说"，Agent 解决"会做"。
  \item \textbf{架构差异}：Agent 有 Agent Loop（推理→行动→观察→再推理的循环），Chatbot 是单轮请求-响应。
\end{itemize}


\medskip
\hrule
\medskip



\section{2. Agent Loop 的完整流程是什么？工程难点在哪里？ \lvchu{}}


\subsection{参考答案要点}

\begin{itemize}[itemsep=1pt, topsep=2pt]
  \item \textbf{完整流程}：
\end{itemize}
\begin{enumerate}[itemsep=1pt, topsep=2pt]
  \item 初始化：加载 System Prompt、可用工具列表、用户请求；
  \item 循环迭代：读取上下文 → LLM 推理下一步（调用工具 or 直接回复）→ 执行工具 → 将结果追加到上下文；
  \item 退出条件：LLM 判断任务完成，不再调用工具时退出；
  \item 安全兜底：设置最大迭代轮次上限（一般 10\textasciitilde\{\}20 轮）或 Token 消耗阈值。
\end{enumerate}
\begin{itemize}[itemsep=1pt, topsep=2pt]
  \item \textbf{工程难点不在 while 循环本身，而在上下文管理}：
  \item 任务越跑越久，上下文越来越长，关键信息被稀释，模型容易跑偏；
  \item 这是 Context Engineering 要解决的核心问题；
  \item 还需要处理工具调用失败、多轮推理跑偏、Token 成本失控等问题。
\end{itemize}


\medskip
\hrule
\medskip



\section{3. ReAct、Plan-and-Execute、Reflection、Multi-Agent 这些范式分别适合什么场景？ \lvzhong{}}


\subsection{参考答案要点}

\begin{itemize}[itemsep=1pt, topsep=2pt]
  \item \textbf{ReAct（Reasoning + Acting）}：
  \item 推理和行动交替进行，走一步看一步；
  \item 适合路径不确定、需要根据证据调整方向的任务（如故障排查、代码库探索）。
  \item \textbf{Plan-and-Execute}：
  \item 先规划整体方案，再逐步执行；
  \item 适合步骤较多但可提前规划的复杂任务，"更像 Workflow 和 Agent 的混合体"；
  \item 相比 ReAct 更可控，Token 消耗更可预测。
  \item \textbf{Reflection}：
  \item 任务完成后进行自我复盘，提取 Meta-Knowledge；
  \item 适合需要持续改进、从失败中学习的场景；
  \item 可做三重验证：真实性验证、交付物验证、数据保真性验证。
  \item \textbf{Multi-Agent}：
  \item 多个 Agent 分工协作（如一个负责规划、一个负责执行、一个负责检查）；
  \item 适合复杂任务可自然拆解的场景；
  \item 存在通信开销、上下文污染、任务一致性问题，生产里不能盲目上。
\end{itemize}


\medskip
\hrule
\medskip



\section{4. Tools 注册时，工具 description 为什么非常关键？ \lvzhong{}}


\subsection{参考答案要点}

\begin{itemize}[itemsep=1pt, topsep=2pt]
  \item \textbf{模型只认结构化描述}：LLM 通过 Function Calling Schema（JSON Schema 格式）理解工具有什么用、什么时候该用、参数怎么填。
  \item \textbf{description 决定了两个核心行为}：
\end{itemize}
\begin{enumerate}[itemsep=1pt, topsep=2pt]
  \item \textbf{是否调用}：模型根据 description 判断当前场景是否适合调用该工具；
  \item \textbf{参数怎么填}：每个参数的 description 告诉模型该字段的含义和合法值范围。
\end{enumerate}
\begin{itemize}[itemsep=1pt, topsep=2pt]
  \item \textbf{最佳实践}：
  \item description 要写清楚工具的适用场景和不适用场景（"什么情况下别调这个"）；
  \item 参数描述要给出格式示例（如 \texttt{"时间范围，格式 HH:MM-HH:MM，比如 09:00-09:30"}）；
  \item 默认值要有明确说明。
  \item \textbf{错误案例}：description 写得模糊，模型会乱调用不相关的工具，造成 Token 浪费和结果偏差。
\end{itemize}


\medskip
\hrule
\medskip



\section{5. Agent 的短期记忆和长期记忆有什么区别？分别怎么落地？ \lvzhong{}}


\subsection{参考答案要点}

\begin{itemize}[itemsep=1pt, topsep=2pt]
  \item \textbf{短期记忆（Session 级）}：
  \item 当前单次会话中的暂存信息：用户提问、模型回复、工具调用中间结果；
  \item 依托 LLM 上下文窗口，会随 Session 结束而消失；
  \item 落地手段：滑动窗口裁剪、上下文摘要压缩、上下文卸载（大结果存外部，Prompt 只放引用）。
  \item \textbf{长期记忆（跨 Session 级）}：
  \item 持久化知识库：用户偏好、历史决策、过往经验；
  \item 通过"写入-检索"机制让新 Session 还能拿到沉淀数据；
  \item 落地手段：向量库 + 语义检索 + Reranker 精排，必要时结合图数据库（Neo4j）做多跳推理。
  \item \textbf{核心差异}：短期是"刚才说过什么"，长期是"过去积累了什么"。两者物理和逻辑上都应分开管理。
\end{itemize}


\medskip
\hrule
\medskip



\section{6. 记忆系统需要解决哪些核心问题？如何避免记忆污染？ \lvgao{}}


\subsection{参考答案要点}

\begin{itemize}[itemsep=1pt, topsep=2pt]
  \item \textbf{核心问题清单}：
\end{itemize}
\begin{enumerate}[itemsep=1pt, topsep=2pt]
  \item 记忆压缩：上下文窗口有限，历史对话需要压缩，但压缩会丢失细节（如精确错误码、关键数值）；
  \item 记忆过期：旧偏好和事实可能已不再适用，需要失效机制；
  \item 记忆冲突：新事实和旧记录矛盾，需要冲突检测与合并；
  \item 记忆检索准确性：向量检索不是精确匹配，语义相近不等于相关。
\end{enumerate}
\begin{itemize}[itemsep=1pt, topsep=2pt]
  \item \textbf{工程策略}：
  \item 权重衰减公式：\texttt{score = relevance × importance × decay(time)}，低分记忆自动淘汰；
  \item 边失效机制（ZEP）：新事实和旧边时间重叠时，标记旧边失效并打时间戳；
  \item LLM 抽取时加校验：JSON Schema 强约束输出、LLM-as-Judge 二次校验、低置信度不写入；
  \item 幂等 Key 设计：基于源消息 ID + 抽取批次 ID，避免重试产生重复记忆。
\end{itemize}


\medskip
\hrule
\medskip



\section{7. 向量记忆和 Markdown 记忆分别适合什么场景？ \lvzhong{}}


\subsection{参考答案要点}

\begin{itemize}[itemsep=1pt, topsep=2pt]
  \item \textbf{向量记忆}：
  \item 适合语义模糊匹配（"用户喜欢简洁的风格"）；
  \item 基于 Embedding + 向量检索，召回率高但不精确；
  \item 适用场景：偏好检索、经验召回、相似情景匹配。
  \item \textbf{Markdown 记忆}：
  \item 适合结构化规则和确定性知识（"该用户的 Java Code Review 优先关注 OOM 风险"）；
  \item 可人工编辑、可 Git 版本管理、可 Code Review；
  \item 适用场景：团队共享知识、Skill 文档、规则库。
  \item \textbf{选型建议}：
  \item 不确定的、需要语义匹配的 → 向量记忆；
  \item 确定的、需要人审核的 → Markdown 记忆；
  \item 二者不是互斥的，可以混合使用。
\end{itemize}


\medskip
\hrule
\medskip



\section{8. 为什么 Agent 场景下只优化 Prompt 不够？Context Engineering 解决什么问题？ \lvgao{}}


\subsection{参考答案要点}

\begin{itemize}[itemsep=1pt, topsep=2pt]
  \item \textbf{Prompt Engineering 的边界}：
  \item 解决"怎么把指令说清楚"的问题，聚焦于单次模型调用的输入质量；
  \item 在简单任务里效果显著，但在长链路 Agent 场景里作用递减。
  \item \textbf{Context Engineering 解决的问题}：
  \item 在合适时机给模型提供正确且必要的信息；
  \item 核心问题：静态规则、动态信息、工具结果、记忆应该如何进入上下文；
  \item 具体手段：Token 降级（JIT 卸载）、上下文隔离（多 Agent 不广播全量历史）、上下文利用率监控（40\% 告警线）。
  \item \textbf{为什么 Prompt 不够}：
  \item Agent 是多轮循环，不是单次调用，信息在循环中持续累积和稀释；
  \item Prompt 调得再好，上下文混乱模型一样跑偏；
  \item 长任务中上下文溢出时，需要 Compaction（压缩）、结构化笔记、Sub-agent 分流等手段，这些都不是 Prompt Engineering 能解决的。
\end{itemize}


\medskip
\hrule
\medskip



\section{9. MCP 协议解决什么问题？它和 Function Calling 有什么区别？ \lvzhong{}}


\subsection{参考答案要点}

\begin{itemize}[itemsep=1pt, topsep=2pt]
  \item \textbf{MCP 解决的问题}：
  \item 工具接入碎片化：每个 LLM 平台各自适配工具，维护成本高；
  \item 提供统一的 Client-Server 协议，让工具以标准化方式被发现和调用；
  \item 被类比为"AI 领域的 USB-C"。
  \item \textbf{三者关系}：MCP Client（使用者，如 Agent）→ MCP Server（提供者，暴露工具）→ Host（运行环境）。
  \item \textbf{MCP 的核心能力}：Tools（执行操作）、Resources（提供数据）、Prompts（模板）。
  \item \textbf{与 Function Calling 的区别}：
  \item \textbf{Function Calling} 解决数据格式问题（OpenAI Schema 描述工具名称/参数/类型）；
  \item \textbf{MCP} 解决通信接入问题（工具如何被发现、连接、调用）；
  \item 二者互补：Function Calling 定义"工具长什么样"，MCP 定义"怎么连上工具"。
  \item \textbf{生产级 MCP Server 安全治理}：权限、审计、限流、脱敏、人工确认都要补上。
\end{itemize}


\medskip
\hrule
\medskip



\section{10. Agent Skills 是什么？Skill 路由怎么做？ \lvgao{}}


\subsection{参考答案要点}

\begin{itemize}[itemsep=1pt, topsep=2pt]
  \item \textbf{Skills 定义}：
  \item 很多任务不是一句 Prompt 能做好，而是需要固定流程、上下文、模板、脚本和校验规则；
  \item Skill 把这些方法封装下来，让 Agent 遇到类似任务时按既定流程执行。
  \item \textbf{Skills 和 Prompt / MCP / Function Calling 的边界}：
  \item Prompt 是"怎么说"，Skill 是"怎么做一套"；
  \item Function Calling 是单次工具调用，Skill 是工具 + 流程 + 上下文 + 校验的组合；
  \item MCP 是工具接入协议，Skill 是更高层的任务能力封装。
  \item \textbf{Skills 为什么要延迟加载}：
  \item 用 \texttt{description} 做路由匹配，命中后再加载 \texttt{SKILL.md} 正文；
  \item 不需要每轮都把规则重新塞进上下文，避免撑爆上下文窗口；
  \item 这也叫"渐进式披露"（Progressive Disclosure）。
  \item \textbf{Skill 路由}：本质是语义匹配问题，和 RAG 相似但目标不同——RAG 检索知识，Skill 检索执行流程。
  \item \textbf{写 SKILL.md 的常见坑}：
  \item 写得像超级 System Prompt，把所有规则塞进去，上下文被撑爆；
  \item 没有写清除边界（什么时候该用、什么时候不该用）；
  \item 没有写失败处理和回滚规则。
\end{itemize}


\medskip
\hrule
\medskip



\section{11. Multi-Agent 协作的主要问题是什么？为什么生产里不能盲目上多 Agent？ \lvzhong{}}


\subsection{参考答案要点}

\begin{itemize}[itemsep=1pt, topsep=2pt]
  \item \textbf{主要问题}：
\end{itemize}
\begin{enumerate}[itemsep=1pt, topsep=2pt]
  \item 上下文污染：子 Agent 拿到不相关的历史信息，越跑越偏；
  \item 通信开销：Agent 之间的消息传递消耗大量 Token；
  \item 任务一致性：多个 Agent 对同一任务的理解可能不一致，产出矛盾；
  \item 局部最优陷阱：每个 Agent 都在自己的局部范围内做最优决策，组合起来全局可能很差；
  \item 调试困难：多 Agent 的决策链路复杂，排查问题困难。
\end{enumerate}
\begin{itemize}[itemsep=1pt, topsep=2pt]
  \item \textbf{生产建议}：
  \item 先用单 Agent + 良好 Harness 把核心链路跑稳；
  \item 只在任务天然可拆解（如"规划→执行→检查"）时才引入多 Agent；
  \item 多 Agent 架构必须做上下文隔离，主 Agent 只给子 Agent 传精简指令；
  \item DAG 编排（如 Coze、Dify）比完全自治的多 Agent 更可控。
\end{itemize}


\medskip
\hrule
\medskip



\section{12. 上下文利用率为什么有一个 40\% 的阈值现象？怎么利用这个现象做工程优化？ \lvzhong{}}


\subsection{参考答案要点}

\begin{itemize}[itemsep=1pt, topsep=2pt]
  \item \textbf{40\% 阈值现象（Smart Zone / Dumb Zone）}：
  \item 来源：Dex Horthy 观察到 168K token 上下文窗口，用到约 40\% 时 Agent 输出质量开始明显下降；
  \item Smart Zone（0\textasciitilde\{\}40\%）：推理聚焦、工具调用准确、代码质量高；
  \item Dumb Zone（>40\%）：幻觉增多、兜圈子、格式混乱、代码变差。
  \item \textbf{原因分析}：
  \item Lost in the Middle：模型对上下文首尾信息利用效率高，中间段利用率明显更低；
  \item 上下文越长，位置偏差越明显，噪声越多。
  \item \textbf{工程优化}：
  \item 把 40\% 当成告警线，超过后触发压缩、分段执行或任务交接；
  \item 用"渐进式披露"和"分层管理"，只让模型看到当前任务相关的上下文；
  \item Context Reset：清空上下文窗口，但通过结构化交接文档保留关键状态；
  \item 在生产环境中监控上下文利用率，被动等待 Agent 跑偏再处理已经晚了。
\end{itemize}


\medskip
\hrule
\medskip



\markboth{RAG 检索增强}{RAG 检索增强}


\section{13. 什么是 RAG？它和传统搜索引擎、微调有什么区别？ \lvchu{}}


\subsection{参考答案要点}

\begin{itemize}[itemsep=1pt, topsep=2pt]
  \item \textbf{RAG 定义}：Retrieval-Augmented Generation，在 LLM 生成回答前，先从外部知识库检索相关内容，注入上下文后再生成。
  \item \textbf{解决什么问题}：
  \item LLM 知识陈旧（训练数据有截止日期）；
  \item 无法访问企业内部私有数据；
  \item 减少幻觉，提供可溯源的引用。
  \item \textbf{与传统搜索引擎的区别}：
  \item 搜索引擎返回链接列表，用户自己判断和阅读；
  \item RAG 检索后由 LLM 自动阅读理解、融合、生成最终答案；
  \item RAG 更省去用户自己查找、比对、总结的步骤。
  \item \textbf{与微调的决策}：
  \item \textbf{RAG}：知识频繁更新、需要可溯源的引用、私有数据不能用于训练；
  \item \textbf{微调}：需要改变模型的行为风格或特定领域推理模式、知识相对稳定；
  \item \textbf{两者结合}：用微调让模型学会"怎么用检索结果"，用 RAG 提供实时知识。
\end{itemize}


\medskip
\hrule
\medskip



\section{14. RAG 系统中 Chunk Size 怎么选？太大和太小分别有什么影响？ \lvchu{}}


\subsection{参考答案要点}

\begin{itemize}[itemsep=1pt, topsep=2pt]
  \item \textbf{Chunk Size 的 Tradeoff}：
  \item \textbf{太小}（如 128 tokens）：语义完整性差，碎片化严重，检索召回率高但准确率低；
  \item \textbf{太大}（如 2048 tokens）：包含过多无关信息，检索精度下降，且会挤占上下文窗口；
  \item \textbf{合适范围}：一般 256\textasciitilde\{\}1024 tokens，需要根据文档类型和数据分布做 A/B 测试。
  \item \textbf{工程经验}：
  \item 技术文档适合中等 chunk（\textasciitilde\{\}512 tokens），保留完整段落语义；
  \item FAQ 类短文本适合小 chunk（\textasciitilde\{\}256 tokens）；
  \item 长文分析类适合大 chunk（\textasciitilde\{\}1024 tokens）或 Sentence Window Retrieval。
  \item \textbf{进阶策略}：
  \item 父子文档（Parent-Child）：小 chunk 检索 + 大 chunk 返回上下文；
  \item 滑动窗口分块：相邻 chunk 之间有重叠，减少语义断裂；
  \item 上下文扩展：检索到的 chunk 自动带上前后文。
\end{itemize}


\medskip
\hrule
\medskip



\section{15. 余弦相似度、内积、欧氏距离分别适合什么 Embedding 模型？ \lvchu{}}


\subsection{参考答案要点}

\begin{itemize}[itemsep=1pt, topsep=2pt]
  \item \textbf{余弦相似度}：
  \item 比较向量方向，对模长不敏感；
  \item 适合大多数经过 L2 归一化的 Embedding 模型（如 OpenAI text-embedding-ada-002）；
  \item 最通用的选择。
  \item \textbf{内积（Dot Product）}：
  \item 同时考虑方向和模长；
  \item 适合未归一化的 Embedding 或需要区分"强相关"和"弱相关"的场景；
  \item 模长通常和信息量正相关。
  \item \textbf{欧氏距离}：
  \item 对模长敏感；
  \item 在文本相似度场景下使用较少，但适合需要空间可解释性的场景；
  \item 配合 L2 归一化后等价于余弦相似度。
  \item \textbf{关键决策点}：先确认 Embedding 模型的输出是否归一化，再选择相似度度量。
\end{itemize}


\medskip
\hrule
\medskip



\section{16. 向量索引 HNSW、IVF、DiskANN 分别适合什么场景？ \lvgao{}}


\subsection{参考答案要点}

\begin{itemize}[itemsep=1pt, topsep=2pt]
  \item \textbf{HNSW（Hierarchical Navigable Small World）}：
  \item 图索引，召回率高（Recall@10 ≥ 0.95），内存占用大；
  \item 适合低并发 + 高精度场景（如 QPS < 50，P99 数十毫秒）；
  \item 元数据过滤率高时 pre-filter 可能破坏图连通性，需要关注。
  \item \textbf{IVF（Inverted File Index）}：
  \item 聚类索引，先粗聚类再精确搜索；
  \item 内存占用小，适合百万级以上数据量；
  \item 召回率略低于 HNSW，但参数调优（nprobe）可以折中。
  \item \textbf{DiskANN}：
  \item 基于 SSD 的图索引，内存占用极小；
  \item 适合十亿级数据量 + 低内存预算的场景；
  \item 延迟比 HNSW 高，但性价比更好。
  \item \textbf{选型维度}：索引类型、元数据过滤方式、多租户隔离、持久化一致性、成本模型。
\end{itemize}


\medskip
\hrule
\medskip



\section{17. RAG 评测为什么必须分检索和生成两段？怎么打分？ \lvzhong{}}


\subsection{参考答案要点}

\begin{itemize}[itemsep=1pt, topsep=2pt]
  \item \textbf{为什么要分段评测}：
  \item 检索和生成是独立环节，问题可能出在任一环节；
  \item 检索不好 → 给 LLM 的上下文本身就是错的，生成再好也没用；
  \item 生成不好 → 检索对了但 LLM 没用对信息；
  \item 混在一起评测无法定位问题根源。
  \item \textbf{检索段评测指标}：
  \item Recall@K：Top-K 结果中召回相关文档的比例；
  \item MRR（Mean Reciprocal Rank）：第一个相关文档的排名倒数均值；
  \item NDCG：考虑排序质量的归一化折损累计增益。
  \item \textbf{生成段评测指标}：
  \item 忠实度（Faithfulness）：生成内容是否基于检索到的文档，有无编造；
  \item 相关性（Relevance）：回答是否切中问题；
  \item 准确性（Correctness）：事实是否正确。
  \item \textbf{评测方法}：
  \item 构建 Golden Set（标准问答对 + 标注相关文档）；
  \item 使用 LLM-as-Judge 自动评测，配合人工抽检校准。
\end{itemize}


\medskip
\hrule
\medskip



\section{18. 长上下文窗口会不会取代 RAG？什么时候用 RAG，什么时候直接用长上下文？ \lvchu{}}


\subsection{参考答案要点}

\begin{itemize}[itemsep=1pt, topsep=2pt]
  \item \textbf{长上下文不取代 RAG 的原因}：
\end{itemize}
\begin{enumerate}[itemsep=1pt, topsep=2pt]
  \item \textbf{成本}：上下文越长，推理成本线性增长；RAG 只注入精选片段，Token 更省；
  \item \textbf{准确性}：Lost in the Middle 问题在长上下文中更严重，中间信息利用率低；
  \item \textbf{信息密度}：RAG 做了检索和排序，给模型的是高密度相关信息，而非海量原始文本；
  \item \textbf{知识更新}：长上下文不能实时更新知识源，RAG 可以随时切换向量库。
\end{enumerate}
\begin{itemize}[itemsep=1pt, topsep=2pt]
  \item \textbf{选型决策}：
  \item 用 \textbf{RAG}：知识量远超上下文窗口、需要精确溯源、知识频繁更新；
  \item 用 \textbf{长上下文}：单文档深度分析（如合同审查）、需要理解全文结构、检索精度难以保证时；
  \item 两者\textbf{结合}：RAG 先粗筛，长上下文再精读，是常见的最佳实践。
\end{itemize}


\medskip
\hrule
\medskip



\section{19. 多模态 RAG 相比纯文本 RAG 多了哪些挑战？ \lvgao{}}


\subsection{参考答案要点}

\begin{itemize}[itemsep=1pt, topsep=2pt]
  \item \textbf{核心挑战}：
\end{itemize}
\begin{enumerate}[itemsep=1pt, topsep=2pt]
  \item Embedding 模型不统一：文本用文本 Embedding，图片用 CLIP 等多模态 Embedding，分属不同向量空间；
  \item Chunk 策略复杂：一个 PPT 页里图表 + 文字 + 表格混排，怎么拆？图文关联怎么保留？
  \item 检索结果融合：用户纯文本 query，怎么召回图片？怎么把图文结果一起排序？
  \item 多模态生成：LLM 能否理解图表并给出文字回答？需要多模态 LLM 支持。
\end{enumerate}
\begin{itemize}[itemsep=1pt, topsep=2pt]
  \item \textbf{工程应对}：
  \item 多向量库并存：文本库 + 图片库，用路由决定查哪个；
  \item 图文关联：图片 Chunk 里带上周围的文字描述或 Caption；
  \item 检索策略：分别检索后做结果融合（Fusion）或交叉排序；
  \item 多模态 LLM：用 GPT-4V / Claude 等多模态模型直接理解图文混合上下文。
\end{itemize}


\medskip
\hrule
\medskip



\section{20. RAG 和 Agent 的记忆系统在技术上有哪些异同？ \lvgao{}}


\subsection{参考答案要点}

\begin{itemize}[itemsep=1pt, topsep=2pt]
  \item \textbf{相同之处}：
  \item 都用向量库和语义检索做相似性匹配；
  \item 都有"写入-检索"两条链路；
  \item 都需要处理检索精度和排序问题（Reranker）。
  \item \textbf{本质区别}：
  \item \textbf{RAG} 挂载的是共享知识源（公司规章、产品文档），非个性化，对不同用户返回同样内容；
  \item \textbf{长期记忆} 管理的是 Agent 与特定用户的交互中沉淀的个性化经验（偏好、习惯、历史决策），高度个性化。
  \item \textbf{配合使用}：
  \item RAG 提供"世界知识"，长期记忆提供"用户画像"；
  \item 检索阶段可以分别召回再融合排序；
  \item 长期记忆里的实体可作为 RAG 检索的 query 扩展；
  \item 用户偏好可作为 RAG 结果的个性化重排信号。
\end{itemize}


\medskip
\hrule
\medskip



\markboth{Agent 工程化}{Agent 工程化}


\section{21. Harness Engineering 是什么？为什么说 Agent = Model + Harness？ \lvzhong{}}


\subsection{参考答案要点}

\begin{itemize}[itemsep=1pt, topsep=2pt]
  \item \textbf{Harness 定义}：
  \item 指模型之外的整套系统：系统提示词、工具调用、文件系统、沙箱环境、编排逻辑、钩子中间件、反馈回路、约束机制；
  \item 模型只提供推理和生成能力，Harness 把状态、工具、反馈、执行环境和安全边界串起来。
  \item \textbf{Agent = Model + Harness 的含义}：
  \item 模型像 CPU，Harness 像操作系统——CPU 再强，OS 天天崩，体验也不会好；
  \item 同一个模型，只改变 Harness（如文件编辑接口的调用方式），编码基准分数可从 6.7\% 跳到 68.3\%。
  \item \textbf{为什么瓶颈经常不在模型}：
  \item 工具接口设计得难用 → 模型调不对；
  \item 反馈回路缺失 → 模型不知道自己错了；
  \item 上下文太杂乱 → 模型输出质量下降；
  \item 调 Prompt 解决不了这些问题。
\end{itemize}


\medskip
\hrule
\medskip



\section{22. Harness 的六层架构分别解决什么问题？ \lvgao{}}


\subsection{参考答案要点}

{\footnotesize
\begin{tabularx}{\textwidth}{|X|X|X|X|}
\hline
\textbf{层级} & \textbf{名称} & \textbf{解决的问题} & \textbf{关键设计} \\
\hline
L1 & 信息边界层 & Agent 该知道什么、不该知道什么 & 定义角色与目标，裁剪无关信息 \\
\hline
L2 & 工具系统层 & Agent 怎么和外部世界交互 & 选择工具、控制调用时机、提炼工具结果 \\
\hline
L3 & 执行编排层 & 多步骤任务怎么串起来 & 理解目标→判断信息→分析→生成→检查 \\
\hline
L4 & 记忆与状态层 & 长任务中间结果怎么管理 & 独立管理任务状态、中间产物和长期记忆 \\
\hline
L5 & 评估与观测层 & Agent 怎么知道自己做对了没有 & 独立于生成过程的验证机制 \\
\hline
L6 & 约束、校验与恢复层 & 出错了怎么办 & 预设规则拦截错误，重试、回滚或降级 \\
\hline
\end{tabularx}
}


\begin{itemize}[itemsep=1pt, topsep=2pt]
  \item \textbf{类比}：L1=岗位说明，L2=办公工具，L3=标准操作流程，L4=项目管理系统，L5=质检，L6=红线规则。
  \item \textbf{实施建议}：不要一上来全搭齐，先做 L1（边界）和 L6（兜底），这两层投入不高但最容易见效。
\end{itemize}


\medskip
\hrule
\medskip



\section{23. 为什么同一个模型换一套工具接口，效果能差 10 倍？ \lvzhong{}}


\subsection{参考答案要点}

\begin{itemize}[itemsep=1pt, topsep=2pt]
  \item \textbf{实验证据}：can.ac 实验，同一个模型，只换了文件编辑接口的调用方式，编码基准分数从 6.7\% 跳到 68.3\%。
  \item \textbf{原因分析}：
\end{itemize}
\begin{enumerate}[itemsep=1pt, topsep=2pt]
  \item 工具返回信息的格式和精度直接影响模型的下一步判断；
  \item 错误信息有没有给出修复方向，决定了模型能否自我纠正；
  \item 工具之间的衔接（如编辑后自动运行 lint，lint 结果自动反馈）形成有效闭环；
  \item 工具接口如果把无关信息也塞给模型，等于在被污染的信息上做推理。
\end{enumerate}
\begin{itemize}[itemsep=1pt, topsep=2pt]
  \item \textbf{model-harness 耦合现象}：
  \item 模型和 Harness 一起调优会过拟合——模型习惯了某套工具逻辑；
  \item 换一个 Harness 表现可能变差；
  \item \textbf{结论}：the best harness for your task is not necessarily the one a model was post-trained with。
\end{itemize}


\medskip
\hrule
\medskip



\section{24. Loop Engineering 是什么？它和 Agent Loop、Workflow Graph 有什么区别？ \lvzhong{}}


\subsection{参考答案要点}

\begin{itemize}[itemsep=1pt, topsep=2pt]
  \item \textbf{Loop Engineering 定义}：
  \item 围绕 Agent 设计可持续运行的反馈循环，让其在明确目标、工具、上下文、验证信号和停止条件下反复行动；
  \item 关注的不再是"一句 Prompt 怎么写"，而是"谁来触发、带哪些材料、跑完拿什么证据判断"。
  \item \textbf{与 Agent Loop（内层循环）的区别}：
  \item Agent Loop：Agent 自己每一轮"推理→行动→观察"的循环；
  \item Engineering Loop（外层循环）：调度系统或人写的流程，唤醒 Agent、分配任务、验证结果、记录状态。
  \item \textbf{与 Workflow Graph 的关系}：
  \item Loop 就是 Graph 里的回边（"生成初稿→审核→不通过继续改"）；
  \item Loop Engineering 把这种循环拓展到跨 Claude Code、CI、GitHub、Issue 系统的日常工作中。
  \item \textbf{核心要素}：触发、目标、上下文、行动、观察、状态、停止条件。
\end{itemize}


\medskip
\hrule
\medskip



\section{25. Workflow、Graph、Loop 三者是什么关系？什么场景用哪种编排方式？ \lvzhong{}}


\subsection{参考答案要点}

\begin{itemize}[itemsep=1pt, topsep=2pt]
  \item \textbf{关系}：
  \item Workflow 是最外层概念，描述"任务怎么做"的流程；
  \item Graph 是 Workflow 的形式化表示（节点 + 边），可以用 DAG（有向无环图）约束执行流程；
  \item Loop 是 Graph 里的回边（条件边），让流程可以循环执行。
  \item \textbf{选型}：
  \item \textbf{纯 Workflow（DAG）}：流程清晰、步骤有限、需要可视化管理和审计（如审批流、数据管道）；
  \item \textbf{Agent Loop（ReAct）}：步骤不确定、需要动态判断的任务（如故障排查）；
  \item \textbf{Plan-and-Execute}：超长流程里夹杂动态子任务，Workflow 和 Agent 的混合体；
  \item \textbf{多 Agent Graph}：任务天然可拆解为并行子任务，由编排层管理 Agent 间消息传递。
  \item \textbf{Graph Loop 的安全设计}：continue 条件、exit 条件、最大轮次、超时、Token 预算、失败降级。
\end{itemize}


\medskip
\hrule
\medskip



\section{26. 如何防止 Agent Loop 陷入死循环？ \lvzhong{}}


\subsection{参考答案要点}

\begin{itemize}[itemsep=1pt, topsep=2pt]
  \item \textbf{多层防护策略}：
\end{itemize}
\begin{enumerate}[itemsep=1pt, topsep=2pt]
  \item 最大迭代轮次上限（硬限制，一般 10\textasciitilde\{\}20 轮）；
  \item Token 消耗上限（防止烧钱）；
  \item 超时限制（防止单轮工具调用挂死）；
  \item 重复检测：如果连续 N 轮调用相同工具且参数相同 → 终止；
  \item 进展检测：如果工具返回结果连续无有效产出 → 终止。
\end{enumerate}
\begin{itemize}[itemsep=1pt, topsep=2pt]
  \item \textbf{工程实现}：
  \item 每轮 loop 记录工具调用和结果摘要；
  \item 设置全局状态变量跟踪任务进展；
  \item 超出限制时，保存部分结果并抛回给用户或转人工。
  \item \textbf{与 Graph Loop 的区别}：Graph Loop 条件更可控（条件边 + 状态机），Agent Loop 更依赖 LLM 自己的判断。
\end{itemize}


\medskip
\hrule
\medskip



\section{27. 大模型网关需要承担哪些核心能力？ \lvgao{}}


\subsection{参考答案要点}

\begin{itemize}[itemsep=1pt, topsep=2pt]
  \item \textbf{核心能力清单}：
\end{itemize}
\begin{enumerate}[itemsep=1pt, topsep=2pt]
  \item \textbf{多模型路由}：根据任务类型、成本、延迟要求路由到不同的模型（如简单任务用小模型，复杂任务用大模型）；
  \item \textbf{Fallback 与容错}：主模型不可用时自动切换到备用模型；
  \item \textbf{限流与流控}：同时看 RPM、TPM、并发数和租户预算四维限流；
  \item \textbf{成本控制}：按 Token 用量、模型类型、租户做成本归因和预算告警；
  \item \textbf{PII 脱敏}：在请求进入模型前做敏感信息检测和脱敏处理；
  \item \textbf{审计与可观测}：记录每次调用的模型、Token 用量、延迟、错误码；
  \item \textbf{协议统一}：屏蔽不同模型 API 的差异，对上游暴露统一接口。
\end{enumerate}
\begin{itemize}[itemsep=1pt, topsep=2pt]
  \item \textbf{为什么不能只按 QPS 限流}：
  \item Token 消耗不均：一次调用可能 10 token 也可能 10K token；
  \item TPM 才是真正的吞吐瓶颈；
  \item 多租户预算隔离需要租户级别的限流维度。
\end{itemize}


\medskip
\hrule
\medskip



\section{28. 模型 Fallback 策略怎么设计？什么时候不能自动降级？ \lvgao{}}


\subsection{参考答案要点}

\begin{itemize}[itemsep=1pt, topsep=2pt]
  \item \textbf{Fallback 层级设计}：
\end{itemize}
\begin{enumerate}[itemsep=1pt, topsep=2pt]
  \item 同厂商同模型重试（瞬态错误）；
  \item 同厂商降级模型（主模型不可用，用同厂小模型）；
  \item 跨厂商切换（切换到另一个云厂商的同等能力模型）；
  \item 静态/缓存回答（所有模型都不可用时的兜底）。
\end{enumerate}
\begin{itemize}[itemsep=1pt, topsep=2pt]
  \item \textbf{什么时候不能自动降级}：
  \item 安全/合规场景（模型回答涉及法律、医疗等高风险决策）；
  \item 结构化输出场景（小模型不保证 JSON Schema 合规）；
  \item 工具调用安全要求高（高风险操作必须有二次确认，不能换模型静默执行）；
  \item 质量敏感场景（降级模型可能导致错误决策的代价高于不可用的损失）。
  \item \textbf{降级策略必须配合}：语义一致性校验、质量阈值、人工审核队列。
\end{itemize}


\medskip
\hrule
\medskip



\section{29. Token 成本怎么归因到租户、用户、功能和 Prompt 版本？ \lvzhong{}}


\subsection{参考答案要点}

\begin{itemize}[itemsep=1pt, topsep=2pt]
  \item \textbf{归因维度设计}：
  \item \textbf{租户维度}：多租户平台必须按租户拆分 Token 用量；
  \item \textbf{用户维度}：租户内部的用户级别的消费统计；
  \item \textbf{功能维度}：不同功能模块（对话、RAG 检索、代码生成、Agent 任务）分别统计；
  \item \textbf{Prompt 版本维度}：不同 Prompt 版本的 Token 效率对比，用于优化。
  \item \textbf{工程实现}：
  \item 每次 API 调用日志记录：\texttt{\{tenant\_id, user\_id, feature, prompt\_version, model, input\_tokens, output\_tokens, cost\}}；
  \item 实时聚合 + 离线批处理两条线：实时用于限流和预算告警，离线用于分析和优化；
  \item 带上异步任务的 Token 消耗（如记忆抽取、检索、反思等后台任务）。
  \item \textbf{常见陷阱}：只算大模型 API 的 Token 成本，忘记算 Embedding、Reranker、向量库查询的算力成本。
\end{itemize}


\medskip
\hrule
\medskip



\section{30. Agent 的 AGENTS.md 应该怎么写？为什么不要写成超级 System Prompt？ \lvchu{}}


\subsection{参考答案要点}

\begin{itemize}[itemsep=1pt, topsep=2pt]
  \item \textbf{AGENTS.md 的定位}：
  \item Agent 每次启动自动加载，是项目级的知识入口；
  \item 应该像目录，而不是把所有规则塞进去。
  \item \textbf{OpenAI 的最佳实践}：
  \item AGENTS.md 大约 100 行，当目录用；
  \item 详细规则放到子文档里按需加载（渐进式披露）；
  \item 每一行对应一个历史失败案例，犯错后更新，形成反馈循环。
  \item \textbf{为什么不能写成超级 System Prompt}：
  \item 上下文窗口有限，全塞进去会进入 Dumb Zone（>40\%利用率质量下降）；
  \item Agent 面对臃肿的规则更容易"迷路"；
  \item 规则混在一起，Agent 不会区分什么场景用什么规则。
  \item \textbf{写法建议}：
  \item 项目基本信息（技术栈、目录结构、关键约定）；
  \item 规则索引（指向子文档）；
  \item 禁止事项（红线规则）；
  \item 常见错误和修复方法。
\end{itemize}


\medskip
\hrule
\medskip



\section{31. 工作流中断后怎么恢复？State 更新策略怎么选？ \lvgao{}}


\subsection{参考答案要点}

\begin{itemize}[itemsep=1pt, topsep=2pt]
  \item \textbf{中断恢复机制}：
\end{itemize}
\begin{enumerate}[itemsep=1pt, topsep=2pt]
  \item 检查点（Checkpoint）：每个关键节点完成后保存 State 快照；
  \item 幂等重放（Replay）：从中断前的最近检查点回放执行；
  \item 人工接管：中断后保存当前状态和上下文，支持人工介入后继续。
\end{enumerate}
\begin{itemize}[itemsep=1pt, topsep=2pt]
  \item \textbf{State 更新策略选择}：
  \item \textbf{Replace（替换）}：适合不变的事实型状态，如用户信息、任务目标；
  \item \textbf{Append（追加）}：适合过程记录型状态，如执行日志、对话历史；
  \item \textbf{Reducer（归约）}：适合需要合并累加的状态，如 Token 消耗累计、进度百分比。
  \item \textbf{实现要点}：
  \item State 要持久化到外部存储（数据库、文件），不能只依赖内存；
  \item 每步更新后原子性写入，避免中断时状态不一致；
  \item State 结构要兼顾人可读（调试）和机器可解析（Schema 约束）。
\end{itemize}


\medskip
\hrule
\medskip



\section{32. AI 应用的可观测性要看哪些指标？ \lvchu{}}


\subsection{参考答案要点}

\begin{itemize}[itemsep=1pt, topsep=2pt]
  \item \textbf{核心指标分类}：
\end{itemize}
\begin{enumerate}[itemsep=1pt, topsep=2pt]
  \item \textbf{延迟指标}：TTFT（首 Token 时间）、端到端延迟、工具调用延迟、检索延迟；
  \item \textbf{质量指标}：任务成功率、幻觉率、用户反馈评分、LLM-as-Judge 评分；
  \item \textbf{成本指标}：Token 消耗、API 调用次数、异步任务开销；
  \item \textbf{可靠性指标}：API 错误率、超时率、Fallback 触发次数、循环异常终止率；
  \item \textbf{安全指标}：PII 检出次数、Prompt 注入尝试次数、高风险操作审计。
\end{enumerate}
\begin{itemize}[itemsep=1pt, topsep=2pt]
  \item \textbf{追踪（Tracing）设计}：
  \item 每次请求生成 traceId，贯穿 Agent Loop 的每一轮；
  \item 记录每轮的输入上下文、LLM 输出、工具调用参数和结果；
  \item 用于事后排查和 Trace 回放评测。
  \item \textbf{告警设计}：
  \item 上下文利用率超过 40\% → 告警；
  \item 任务成功率连续下降 → 告警（可能是模型/Prompt/检索出现变化）；
  \item Token 消耗异常飙升 → 告警。
\end{itemize}


\medskip
\hrule
\medskip



\markboth{系统设计}{系统设计}


\section{33. 请设计一个生产级 AI 应用的整体架构，从入口到模型返回的完整链路是什么？ \lvgao{}}


\subsection{参考答案要点}

\begin{itemize}[itemsep=1pt, topsep=2pt]
  \item \textbf{六层架构}（自外向内）：
\end{itemize}

{\footnotesize
\begin{tabularx}{\textwidth}{|X|X|}
\hline
\textbf{层级} & \textbf{职责} \\
\hline
入口层 & API Gateway、认证鉴权、限流、请求路由 \\
\hline
编排层 & 理解用户意图 → 判断是否需要 RAG/Agent/Tool → 编排执行计划 \\
\hline
Prompt/Context 层 & 组装 System Prompt、注入 RAG 结果、注入记忆、Token 预算控制 \\
\hline
RAG/Memory/Tool 层 & 检索知识、召回记忆、执行工具调用、结果提炼 \\
\hline
模型网关层 & 多模型路由、Fallback、重试、限流、PII 脱敏、审计 \\
\hline
评测观测层 & 日志/Trace 采集、质量评测、成本归因、告警 \\
\hline
\end{tabularx}
}


\begin{itemize}[itemsep=1pt, topsep=2pt]
  \item \textbf{请求链路}：
\end{itemize}
\begin{enumerate}[itemsep=1pt, topsep=2pt]
  \item 用户请求进入 API Gateway → 鉴权 + 限流；
  \item 编排层分析意图，决定走简单问答/RAG/Agent Loop；
  \item Context 层注入相关记忆 + RAG 结果；
  \item 如需工具调用，Agent Loop 执行推理→调用→观察循环；
  \item 模型网关选模型、发请求、处理重试和降级；
  \item 结果经过安全校验后返回用户，同时写入日志和 Trace。
\end{enumerate}


\medskip
\hrule
\medskip



\section{34. 同步、流式（SSE）、异步三种模式怎么选？ \lvchu{}}


\subsection{参考答案要点}

{\footnotesize
\begin{tabularx}{\textwidth}{|X|X|X|}
\hline
\textbf{模式} & \textbf{适合场景} & \textbf{优缺点} \\
\hline
同步（Request-Response） & 短回复、低延迟要求 & 实现简单，用户等待时间长 \\
\hline
流式（SSE） & 对话式 AI、逐字生成 & 显著改善体感（TTFT 低），不能减少总耗时，需处理断连重连 \\
\hline
异步（任务队列） & 长任务（Agent、报告生成） & 用户不阻塞，需轮询或 Webhook 通知结果，复杂度高 \\
\hline
\end{tabularx}
}


\begin{itemize}[itemsep=1pt, topsep=2pt]
  \item \textbf{选型建议}：
  \item 简单问答 → SSE 流式输出；
  \item Agent 长任务 → 异步模式，任务提交后返回 taskId，结果通过 WebSocket 推送或轮询获取；
  \item WebSocket vs SSE：WebSocket 双向通信，适合语音 Agent 等需要双向推送的场景；SSE 更轻量，适合纯单向流式文本。
\end{itemize}


\medskip
\hrule
\medskip



\section{35. 为什么大模型调用限流要同时看 RPM、TPM、并发数和租户预算？ \lvzhong{}}


\subsection{参考答案要点}

\begin{itemize}[itemsep=1pt, topsep=2pt]
  \item \textbf{为什么不能只按 RPM 限流}：
  \item RPM（Request Per Minute）只控制了请求频率；
  \item 一次请求可能消耗 10 token，也可能消耗 10 万 token；
  \item 短请求高频 vs 长请求低频，RPM 相同的 TPM 可能差 100 倍。
  \item \textbf{四维限流体系}：
  \item \textbf{RPM}：控制 API 调用频率（防止连接数爆炸）；
  \item \textbf{TPM}：控制真正的吞吐瓶颈（Token Per Minute）；
  \item \textbf{并发数}：控制同时正在处理的请求数（保护后端服务）；
  \item \textbf{租户预算}：控制每个租户的成本上限，防止单租户超预算。
  \item \textbf{工程实现}：
  \item 分层限流：全局 → 租户 → 用户 → 功能；
  \item Token 计算器在每次请求前后原子性更新；
  \item 超限时返回 429 + Retry-After，而非直接丢弃。
\end{itemize}


\medskip
\hrule
\medskip



\section{36. AI 应用如何做超时、重试、限流、熔断和降级？ \lvgao{}}


\subsection{参考答案要点}

\begin{itemize}[itemsep=1pt, topsep=2pt]
  \item \textbf{超时策略}：
  \item 分级超时：LLM API 调用 30s，工具调用按工具类型设不同超时；
  \item 流式请求用首 Token 超时（如 5s 没收到首 Token 就 Fallback）。
  \item \textbf{重试策略}：
  \item 可重试：网络错误、429 限流、5xx 临时故障；
  \item 不可重试：4xx 参数错误、402 余额不足、内容安全拦截；
  \item 退避策略：指数退避 + 随机抖动（jitter）。
  \item \textbf{熔断}：
  \item 错误率超过阈值 → 熔断器打开，快速失败；
  \item 半开状态：熔断后定期试探性调用，成功则恢复正常。
  \item \textbf{降级}：
  \item 模型降级：大模型 → 小模型 → 静态规则回答；
  \item 功能降级：去掉 RAG 检索直接用模型知识回答；
  \item 结果降级：无模型可用时返回缓存或预设兜底信息。
\end{itemize}


\medskip
\hrule
\medskip



\section{37. 从 Prompt Demo 到生产级架构，最大的差距在哪？要补齐哪些能力？ \lvgao{}}


\subsection{参考答案要点}

\begin{itemize}[itemsep=1pt, topsep=2pt]
  \item \textbf{最大差距}：Demo 只验证"能不能跑"，生产需要"跑得稳、跑得对、跑得划算、跑得安全"。
  \item \textbf{要补齐的能力（按优先级）}：
\end{itemize}

{\footnotesize
\begin{tabularx}{\textwidth}{|X|X|}
\hline
\textbf{阶段} & \textbf{补齐内容} \\
\hline
P0（必须） & 错误处理与重试、Prompt 版本管理、基本日志 \\
\hline
P1（重要） & 模型网关（路由+Fallback+限流）、RAG 检索优化、评测体系（Golden Set + LLM-as-Judge） \\
\hline
P2（增强） & 记忆系统、工具调用安全治理、成本归因、A/B 实验、可观测（Trace + Metrics） \\
\hline
\end{tabularx}
}


\begin{itemize}[itemsep=1pt, topsep=2pt]
  \item \textbf{架构演进路径}：
  \item Prompt Demo → 加上错误处理 → 引入 RAG → 模型网关 → Agent Loop → Harness → 完整评测体系；
  \item 核心原则：MVP 先跑通，再逐步补齐非功能需求。
\end{itemize}


\medskip
\hrule
\medskip



\section{38. Prompt 注入攻击在系统设计层面怎么防？ \lvgao{}}


\subsection{参考答案要点}

\begin{itemize}[itemsep=1pt, topsep=2pt]
  \item \textbf{多层防御策略}：
\end{itemize}
\begin{enumerate}[itemsep=1pt, topsep=2pt]
  \item \textbf{输入层}：检测和过滤可疑的 Prompt 注入模式（如"忽略之前的指令""你现在是 DAN"）；
  \item \textbf{权限层}：工具调用最小权限，高风险操作必须二次人工确认；
  \item \textbf{隔离层}：不可信输入和执行环境隔离（沙箱），Agent 代码执行在受限容器中；
  \item \textbf{上下文层}：System Prompt 使用特殊分隔符，明确区分"系统指令"和"用户输入"；
  \item \textbf{审计层}：记录所有 Prompt 注入检测日志和高风险操作审计。
\end{enumerate}
\begin{itemize}[itemsep=1pt, topsep=2pt]
  \item \textbf{常见注入模式}：
  \item 直接覆盖：\texttt{Ignore all previous instructions and...}
  \item 间接注入：通过检索到的文档或网页内容注入（RAG 场景）→ 需要对检索结果做安全过滤；
  \item 工具结果注入：攻击者控制的 API 返回恶意指令 → 工具结果也需要安全校验。
  \item \textbf{局限}：Prompt 注入没有完美防御方案，需要在安全性和可用性之间做权衡。
\end{itemize}


\medskip
\hrule
\medskip



\section{39. 如何设计一个实时语音 Agent 系统？ \lvgao{}}


\subsection{参考答案要点}

\begin{itemize}[itemsep=1pt, topsep=2pt]
  \item \textbf{核心组件}：
  \item \textbf{ASR}：语音转文字（Whisper、Deepgram 等）；
  \item \textbf{LLM}：文本推理和决策（GPT-4o、Claude 等）；
  \item \textbf{TTS}：文字转语音（ElevenLabs、Edge TTS 等）；
  \item \textbf{VAD}：语音活动检测，判断用户是否在说话。
  \item \textbf{端到端延迟来源}：ASR 延迟 + LLM 首 Token 延迟 + TTS 首音频延迟 + 网络传输延迟。
  \item \textbf{状态管理}：
  \item listening（监听中）→ thinking（推理中）→ speaking（播报中）→ interrupted（被打断）；
  \item 用户打断时：取消 TTS 播放 → 取消 LLM 生成 → 更新上下文为"用户打断了上次回答"；
  \item 使用 WebSocket 维护双向全双工通道。
  \item \textbf{架构选择}：
  \item \textbf{云端 API}：延迟在 500ms-2s，成本可控，适合大多数场景；
  \item \textbf{端云混合}：ASR 和 TTS 在本地，LLM 在云端，延迟可低至 200-800ms；
  \item \textbf{Speech-to-Speech API}：跳过 ASR/TTS 文本中间层，延迟最低但灵活度也最低。
  \item \textbf{可观测指标}：端到端延迟、打断响应延迟、ASR 准确率、对话成功率、用户满意度。
\end{itemize}


\medskip
\hrule
\medskip



\section{40. 出现模型输出事故后，如何通过 Trace 回放定位问题？ \lvgao{}}


\subsection{参考答案要点}

\begin{itemize}[itemsep=1pt, topsep=2pt]
  \item \textbf{Trace 应该记录的信息}：
  \item 每轮 Agent Loop：输入上下文、LLM 请求参数（model、temperature、messages）、LLM 原始输出；
  \item 每次工具调用：工具名、入参、返回结果、耗时；
  \item 每次检索：query、检索结果（文档 ID、chunk、相似度分数）；
  \item 最终输出：经过后处理的返回给用户的回答。
  \item \textbf{回放排查步骤}：
\end{itemize}
\begin{enumerate}[itemsep=1pt, topsep=2pt]
  \item 从用户反馈的 traceId 找到本次请求的完整链路；
  \item 检查 LLM 输出：幻觉？格式错误？工具调用参数填错？
  \item 检查检索结果：检索到的文档是否相关？是否包含了错误信息？
  \item 检查上下文：上下文是否已超过 40\% 阈值？是否有之前工具结果污染？
  \item 定位根因后，用相同输入回放测试，验证修复方案；
  \item 将本次事故加入 Golden Set 作为回归用例。
\end{enumerate}
\begin{itemize}[itemsep=1pt, topsep=2pt]
  \item \textbf{系统要求}：Trace 必须能覆盖全链路，不能只有 LLM 调用的日志。
\end{itemize}


\medskip
\hrule
\medskip



\section{41. 如何构建 Golden Set？冷启动阶段没有生产日志怎么办？ \lvzhong{}}


\subsection{参考答案要点}

\begin{itemize}[itemsep=1pt, topsep=2pt]
  \item \textbf{Golden Set 的覆盖维度}：
  \item 正常路径（Happy Path）：最常见的用户场景；
  \item 边缘场景（Edge Cases）：边界输入、模糊表达、长尾问题；
  \item 对抗样本（Adversarial）：试图注入、越权、误导的输入；
  \item 高权重失败（High-Stakes）：答错后果严重的场景。
  \item \textbf{冷启动策略}：
\end{itemize}
\begin{enumerate}[itemsep=1pt, topsep=2pt]
  \item 人工构造：产品经理 + 领域专家编写 50\textasciitilde\{\}100 条核心问答对；
  \item 文档采样：从知识库文档中抽取代表性段落，构造对应的问句；
  \item 历史数据迁移：如果已有旧的客服记录或 FAQ，清洗后使用；
  \item 红队测试：安全团队构造对抗样本，测试边界行为。
\end{enumerate}
\begin{itemize}[itemsep=1pt, topsep=2pt]
  \item \textbf{持续迭代}：
  \item 将生产中发现的新问题持续加入 Golden Set；
  \item 定期人工抽检 LLM-as-Judge 评分，校准评测模型偏差；
  \item Golden Set 要版本管理，A/B 实验对照。
\end{itemize}


\medskip
\hrule
\medskip



\section{42. LLM-as-Judge 有哪些主要偏差？怎么校准？ \lvgao{}}


\subsection{参考答案要点}

\begin{itemize}[itemsep=1pt, topsep=2pt]
  \item \textbf{主要偏差类型}：
\end{itemize}
\begin{enumerate}[itemsep=1pt, topsep=2pt]
  \item \textbf{位置偏差}：更倾向排在前面的候选答案；
  \item \textbf{冗长偏好}：更容易给长答案高分，即使信息密度不高；
  \item \textbf{自信偏好}：语气更自信的回答评分更高，即使事实错误；
  \item \textbf{自我增强偏差}：更偏向和自身训练数据风格一致的答案；
  \item \textbf{顺序效应}：评测顺序影响打分，前面的更容易得高分或低分。
\end{enumerate}
\begin{itemize}[itemsep=1pt, topsep=2pt]
  \item \textbf{校准方法}：
  \item 答案位置随机交换后评两次，取平均；
  \item 评测 rubric 要具体，不用笼统的"好/不好"，而是按维度打分（忠实度、相关性、准确性）；
  \item 人工抽检 10\%\textasciitilde\{\}20\% 的评测结果，计算人工与 LLM-as-Judge 的一致性；
  \item 不一致超过阈值时，调整评测 Prompt 或 rubric；
  \item 多 LLM 交叉评测，取投票结果。
  \item \textbf{CI 中的平衡}：评测集大小和评测频率要控制，避免 CI 时间过长和 API 成本过高。
\end{itemize}


\medskip
\hrule
\medskip



\section{43. Agent 评测为什么比普通 RAG 和问答评测复杂得多？ \lvzhong{}}


\subsection{参考答案要点}

\begin{itemize}[itemsep=1pt, topsep=2pt]
  \item \textbf{复杂性来源}：
\end{itemize}
\begin{enumerate}[itemsep=1pt, topsep=2pt]
  \item \textbf{多步推理}：Agent 不是单次问答，而是一个序列决策过程，每一步都可能出错；
  \item \textbf{工具调用链}：不仅评最终结果，还要评工具选择是否正确、调用参数是否合理、工具结果是否被正确利用；
  \item \textbf{非确定性}：Agent 执行路径随采样参数、上下文变化而不同，同一输入可能产生不同执行路径；
  \item \textbf{中间状态评估}：需要评估 Agent 是否陷入了死循环、是否做出了不必要的工具调用；
  \item \textbf{安全维度}：Agent 能执行代码、调 API、读写文件，评测要覆盖安全边界。
\end{enumerate}
\begin{itemize}[itemsep=1pt, topsep=2pt]
  \item \textbf{评测维度扩展}：
  \item 除了最终答案质量，还要评：规划合理性、工具选择准确率、Loop 效率（轮次/Token 消耗）、安全合规得分。
  \item \textbf{评测方法}：
  \item 离线评测：用 Golden Set 跑批，评最终结果和关键中间步骤；
  \item Trace 回放：用生产 Trace 重新执行，看修复后结果是否一致；
  \item 线上灰度：小流量对比新旧 Agent，通过用户反馈和业务指标评判。
\end{itemize}


\medskip
\hrule
\medskip



\section{44. RAG、Agent、结构化输出的评测指标为什么不能混用一套？ \lvchu{}}


\subsection{参考答案要点}

\begin{itemize}[itemsep=1pt, topsep=2pt]
  \item \textbf{本质差异}：
  \item \textbf{RAG} 评的是"检索对没对 + 生成对没对"，关注忠实度、召回率；
  \item \textbf{Agent} 评的是"决策链对不对 + 执行到位没有"，关注任务成功率、工具调用准确性、执行效率；
  \item \textbf{结构化输出} 评的是"格式合不合 Schema + 内容对不对"，关注 Schema 合规率、字段级准确率。
  \item \textbf{为什么不能混用}：
  \item 打分维度和权重完全不同：RAG 重视检索精度，Agent 重视规划合理性，结构化输出重视格式合规率；
  \item 失败模式不同：RAG 失败通常是检索不到文档，Agent 失败可能是死循环或工具调用错误，结构化输出失败是格式解析失败；
  \item 使用同一套 rubric 会导致打分不准确，无法定位具体问题。
  \item \textbf{建议}：三类场景分别建立评测集 + 评测 Prompt + 评测指标，再在 CI 流水线中串行执行。
\end{itemize}


\medskip
\hrule
\medskip



\markboth{LLM 基础与结构化输出}{LLM 基础与结构化输出}


\section{45. Token 是什么？为什么中文、英文、代码消耗的 Token 不一样？ \lvchu{}}


\subsection{参考答案要点}

\begin{itemize}[itemsep=1pt, topsep=2pt]
  \item \textbf{Token 定义}：
  \item LLM 处理文本的最小单元，不等于"一个字"或"一个词"；
  \item 取决于分词器（Tokenizer），常见方式：BPE（Byte Pair Encoding）。
  \item \textbf{为什么消耗不同}：
  \item 英文：一个常见单词通常 1\textasciitilde\{\}2 个 Token（如 "hello" = 1 token）；
  \item 中文：一个汉字通常 1\textasciitilde\{\}3 个 Token（如 "你好" 可能 = 2\textasciitilde\{\}6 tokens）；
  \item 代码：符号密集（括号、换行），每个符号通常 1 个 Token。
  \item \textbf{工程影响}：
  \item 中英文混合场景 Token 估算不准；
  \item 代码生成 Token 消耗比看起来大（缩进空格也是 Token）。
\end{itemize}


\medskip
\hrule
\medskip



\section{46. Temperature、Top-P、Top-K 分别控制什么？生产环境怎么设置更稳？ \lvchu{}}


\subsection{参考答案要点}

{\footnotesize
\begin{tabularx}{\textwidth}{|X|X|X|}
\hline
\textbf{参数} & \textbf{控制什么} & \textbf{生产建议} \\
\hline
Temperature & 概率分布的"陡峭度"，越低越确定 & 事实性任务 0\textasciitilde\{\}0.3，创意任务 0.7\textasciitilde\{\}1.0 \\
\hline
Top-P & 只考虑累积概率达到 P 的 token & 0.9\textasciitilde\{\}0.95，过滤尾部低概率 token \\
\hline
Top-K & 只考虑概率最高的 K 个 token & 40\textasciitilde\{\}60，适合需要多样性但不过分发散 \\
\hline
\end{tabularx}
}


\begin{itemize}[itemsep=1pt, topsep=2pt]
  \item \textbf{为什么 Temperature=0 输出仍可能不完全一致}：
  \item 浮点运算精度差异；
  \item 不同硬件/框架的矩阵运算顺序可能不同；
  \item 真正需要确定性应使用 seed 参数固定随机数生成器。
  \item \textbf{生产建议}：
  \item 事实性任务（RAG 回答、代码生成、数据提取）→ Temperature=0 配合固定 seed；
  \item 开放式对话 → Temperature=0.7 + Top-P=0.9；
  \item 不要三个参数同时调，先固定两个，调第三个看效果。
\end{itemize}


\medskip
\hrule
\medskip



\section{47. Streaming 为什么能改善用户体验？它能减少总耗时和 Token 成本吗？ \lvchu{}}


\subsection{参考答案要点}

\begin{itemize}[itemsep=1pt, topsep=2pt]
  \item \textbf{为什么能改善体验}：
  \item 降低 TTFT（Time to First Token）：用户不用等完整回答，首字秒出；
  \item 给予心理反馈：用户知道系统"在干活"，不是卡住了。
  \item \textbf{不能减少什么}：
  \item 总耗时不会减少（甚至可能稍长，因为多次网络传输开销）；
  \item Token 成本不变（模型还是要生成完整回答）。
  \item \textbf{技术选型}：
  \item \textbf{SSE}：轻量、单向、HTTP 标准，适合普通 Chat 流式输出；
  \item \textbf{WebSocket}：全双工、适合语音 Agent 和需要双向通信的场景；
  \item \textbf{HTTP Chunked Transfer}：最底层方案，兼容性好但需要手动解析。
  \item \textbf{工程注意}：
  \item 断连重连：SSE 断连后需要重建上下文，可能造成 Token 浪费；
  \item 取消流：用户中途打断，需要向 LLM API 发送取消信号停止计费。
\end{itemize}


\medskip
\hrule
\medskip



\section{48. JSON Mode 和 Structured Outputs 有什么区别？ \lvzhong{}}


\subsection{参考答案要点}

\begin{itemize}[itemsep=1pt, topsep=2pt]
  \item \textbf{JSON Mode}：
  \item 只保证输出是合法的 JSON 格式；
  \item 不保证 JSON 的结构和字段符合约定 Schema；
  \item 依赖 Prompt 中描述期望的字段，可靠性有限。
  \item \textbf{Structured Outputs（JSON Schema）}：
  \item 在 API 层面强约束输出必须符合给定的 JSON Schema；
  \item 会修改模型采样逻辑，跳过不符合 Schema 的 token；
  \item 可靠性远高于 JSON Mode，几乎 100\% 格式合规。
  \item \textbf{为什么只写"请返回 JSON"不可靠}：
  \item LLM 本质是概率采样，无法保证每次都生成合法 JSON；
  \item 嵌套结构或复杂字段时更容易出错；
  \item 错误率随输出长度增加而上升。
  \item \textbf{结构化输出失败怎么处理}：
  \item 重试（带原输出和错误信息）；
  \item 用正则/解析器尝试修复（如补全括号）；
  \item 降级为非结构化回答（用户体验有损但系统可用）。
\end{itemize}


\medskip
\hrule
\medskip



\section{49. Function Calling 的完整链路是什么？和 MCP Agent Skill 是什么关系？ \lvzhong{}}


\subsection{参考答案要点}

\begin{itemize}[itemsep=1pt, topsep=2pt]
  \item \textbf{Function Calling 完整链路}：
\end{itemize}
\begin{enumerate}[itemsep=1pt, topsep=2pt]
  \item 组装请求：System Prompt + 工具列表（JSON Schema 描述）+ 用户消息；
  \item LLM 推理：决定是否需要调用工具，选择工具，填写参数 JSON；
  \item 框架解析 LLM 输出中的 tool\_call 字段；
  \item 执行工具函数，获取返回结果；
  \item 将工具结果追加到消息列表（role: tool）；
  \item 再次调用 LLM，让其基于工具结果继续推理或生成最终回答。
\end{enumerate}
\begin{itemize}[itemsep=1pt, topsep=2pt]
  \item \textbf{三者关系}：
\end{itemize}

{\footnotesize
\begin{tabularx}{\textwidth}{|X|X|X|X|}
\hline
\textbf{} & \textbf{Function Calling} & \textbf{MCP} & \textbf{Agent Skill} \\
\hline
层级 & 数据格式层 & 通信协议层 & 任务能力封装层 \\
\hline
解决的问题 & 工具"长什么样" & 工具"怎么连" & 任务"怎么做一套" \\
\hline
粒度 & 单次工具调用 & 工具发现+调用 & 工具+流程+上下文+校验 \\
\hline
\end{tabularx}
}


\begin{itemize}[itemsep=1pt, topsep=2pt]
  \item \textbf{一句话概括}：Function Calling 定义方法签名，MCP 提供远程调用通道，Skill 把多个方法调用封装成可复用的任务模板。
\end{itemize}


\medskip
\hrule
\medskip



\section{50. 工具调用为什么必须做安全治理？安全边界在哪里？ \lvgao{}}


\subsection{参考答案要点}

\begin{itemize}[itemsep=1pt, topsep=2pt]
  \item \textbf{安全风险}：
\end{itemize}
\begin{enumerate}[itemsep=1pt, topsep=2pt]
  \item LLM 可能被诱导调用不该调的工具（Prompt 注入）；
  \item LLM 填写的参数可能不安全（文件路径穿越、SQL 注入、Shell 注入）；
  \item 工具调用结果可能包含敏感信息被带入后续上下文；
  \item 无人值守 Agent 可能反复调用昂贵接口（费用失控）；
  \item Agent 可能修改不应修改的代码或配置。
\end{enumerate}
\begin{itemize}[itemsep=1pt, topsep=2pt]
  \item \textbf{安全边界设计}：
  \item \textbf{权限最小化}：每个 Agent 只能调用必需的、最小权限的工具集；
  \item \textbf{参数校验}：工具函数内部做严格的参数校验（白名单、正则、范围检查）；
  \item \textbf{二次确认}：高风险操作（删除、发布、数据库变更）必须人工确认；
  \item \textbf{沙箱隔离}：代码执行在容器/沙箱中，限制文件系统和网络访问；
  \item \textbf{审计日志}：记录每次工具调用的入参、结果、耗时和触发上下文；
  \item \textbf{限流}：每个工具设置独立的调用频率和 Token 消耗上限。
\end{itemize}


\medskip
\hrule
\medskip



\section{51. 哪些大模型 API 错误可以重试？哪些不能重试？为什么必须做幂等？ \lvchu{}}


\subsection{参考答案要点}

\begin{itemize}[itemsep=1pt, topsep=2pt]
  \item \textbf{可重试错误}：
  \item 429 Rate Limit：等待 Retry-After 后重试；
  \item 5xx Server Error：瞬时故障，指数退避重试（最多 3 次）；
  \item 网络超时/连接重置：重试。
  \item \textbf{不可重试错误}：
  \item 401/403 鉴权失败：重试不会变好；
  \item 400 参数错误：需修复请求后再试；
  \item 402 余额不足：需充值；
  \item 内容安全拦截（如 moderation flag）：换内容。
  \item \textbf{为什么必须做幂等}：
  \item 网络波动可能导致请求已处理但响应丢失；
  \item 重试时服务端可能重复处理（重复扣费、重复执行副作用）；
  \item 实现方式：请求带上 idempotency key，服务端做去重。
\end{itemize}


\medskip
\hrule
\medskip



\section{52. 大模型为什么会产生幻觉？常见缓解方案有哪些？ \lvchu{}}


\subsection{参考答案要点}

\begin{itemize}[itemsep=1pt, topsep=2pt]
  \item \textbf{幻觉的根本原因}：
\end{itemize}
\begin{enumerate}[itemsep=1pt, topsep=2pt]
  \item LLM 本质是概率模型，不是知识库——它在预测下一个 token，不是检索事实；
  \item 训练数据中的错误和不一致会被模型学习；
  \item 模型不知道"自己不知道"，强答不擅长的问题；
  \item 上下文信息不足或矛盾时，模型倾向自己"补全"。
\end{enumerate}
\begin{itemize}[itemsep=1pt, topsep=2pt]
  \item \textbf{缓解方案（从上到下优先级）}：
\end{itemize}
\begin{enumerate}[itemsep=1pt, topsep=2pt]
  \item \textbf{RAG}：给模型注入权威的、可溯源的知识片段，最有效的缓释手段；
  \item \textbf{结构化约束}：用 JSON Schema / Structured Outputs 约束输出格式，减少自由发挥空间；
  \item \textbf{Prompt 引导}：明确告诉模型"不知道就说不知道"，禁止编造；
  \item \textbf{Truthfulness 校验}：用 LLM-as-Judge 对关键回答做事实核查；
  \item \textbf{人工审核}：高风险场景加入工 Review 环节。
\end{enumerate}
\begin{itemize}[itemsep=1pt, topsep=2pt]
  \item \textbf{不能完全消灭幻觉}：RAG 减少幻觉但不能根除，LLM 仍可能在推理步骤里生成错误判断。
\end{itemize}


\medskip
\hrule
\medskip



\section{53. Token 预算怎么估算？输入、输出、历史消息、RAG 证据如何取舍？ \lvzhong{}}


\subsection{参考答案要点}

\begin{itemize}[itemsep=1pt, topsep=2pt]
  \item \textbf{Token 预算公式}：
\end{itemize}

\begin{lstlisting}
  上下文窗口剩余 = 模型最大上下文 - System Prompt Token - 历史对话 Token - RAG 证据 Token - 预留输出 Token
\end{lstlisting}

\begin{itemize}[itemsep=1pt, topsep=2pt]
  \item \textbf{各部分取舍策略}：
  \item \textbf{System Prompt}：尽量精简（200\textasciitilde\{\}500 tokens），长规则放 SKILL.md 按需加载；
  \item \textbf{历史对话}：滑动窗口保留最近 N 轮，更老的做摘要压缩；
  \item \textbf{RAG 证据}：Top-3 到 Top-5 个 chunk，每个 256\textasciitilde\{\}512 tokens，总量控制 1500\textasciitilde\{\}2500 tokens；
  \item \textbf{预留输出}：根据任务类型预留 500\textasciitilde\{\}2000 tokens 给模型回答。
  \item \textbf{预算超支处理}：
  \item 先削减 RAG 证据（减少 Top-K 或缩短 chunk）；
  \item 再压缩历史对话（摘要 / 滑动窗口）;
  \item 极端情况：切换到更大的上下文窗口模型。
  \item \textbf{实时监控}：每次调用前计算 Token 使用率，超过 40\% 触发告警。
\end{itemize}


\medskip
\hrule
\medskip



\section{54. AI 应用调用日志里至少要记录哪些字段？为什么？ \lvchu{}}


\subsection{参考答案要点}

\begin{itemize}[itemsep=1pt, topsep=2pt]
  \item \textbf{最小必录字段}：
  \item 基础：\texttt{timestamp}, \texttt{trace\_id}, \texttt{user\_id}, \texttt{tenant\_id}, \texttt{feature};
  \item 请求：\texttt{model}, \texttt{temperature}, \texttt{input\_tokens}, \texttt{system\_prompt\_hash}, \texttt{rag\_query}；
  \item 响应：\texttt{output\_tokens}, \texttt{total\_tokens}, \texttt{latency\_ms}, \texttt{ttft\_ms}, \texttt{finish\_reason}；
  \item 工具：\texttt{tool\_calls[]}（name, params, result\_summary, duration）；
  \item 错误：\texttt{error\_code}, \texttt{error\_message}, \texttt{retry\_count}；
  \item 成本：\texttt{cost\_usd}, \texttt{prompt\_version}。
  \item \textbf{为什么记录这些}：
  \item 成本归因：按 user/tenant/feature/prompt\_version 聚合；
  \item 质量分析：finish\_reason 区分正常完成/截断/内容过滤；
  \item 故障排查：trace\_id 串联全链路，error\_code 定位故障点；
  \item 优化基础：input/output tokens 对比优化效果。
  \item \textbf{日志设计原则}：
  \item 结构化日志（JSON），不要用自由文本；
  \item 敏感字段脱敏（用户消息内容 hash 后存储）。
\end{itemize}


\medskip
\hrule
\medskip



\section{55. 如果 LLM-as-Judge 和人工评测结果不一致，应该怎么处理？ \lvzhong{}}


\subsection{参考答案要点}

\begin{itemize}[itemsep=1pt, topsep=2pt]
  \item \textbf{不一致的常见原因}：
\end{itemize}
\begin{enumerate}[itemsep=1pt, topsep=2pt]
  \item 评测 rubric 不够具体，"好"的标准模糊；
  \item LLM-as-Judge 有系统性偏差（位置偏差、冗长偏好等）；
  \item 人工评测者有自己的主观偏好；
  \item 评测样本数量不足，结论不具有统计意义。
\end{enumerate}
\begin{itemize}[itemsep=1pt, topsep=2pt]
  \item \textbf{处理流程}：
\end{itemize}
\begin{enumerate}[itemsep=1pt, topsep=2pt]
  \item \textbf{抽样对比}：从不一致的样本中随机抽 20\textasciitilde\{\}30 条，由多位人工评测者独立打分；
  \item \textbf{找出分歧模式}：LLM-as-Judge 是否系统性地高估/低估某一类回答？
  \item \textbf{调整 rubric}：细化评测维度（如把"回答质量"拆成忠实度、相关性、准确性三个子维度）；
  \item \textbf{校准评测 Prompt}：在 LLM-as-Judge 的 prompt 中针对发现的偏差加反例约束；
  \item \textbf{计算 IAA}：计算 LLM-as-Judge 与人工评测的 Inter-Annotator Agreement，目标 > 0.7；
  \item \textbf{持续迭代}：每次发现新的偏差类型后，更新评测 Prompt 和 rubric。
\end{enumerate}
\begin{itemize}[itemsep=1pt, topsep=2pt]
  \item \textbf{兜底方案}：高权重决策（如安全、合规）最终以人工评测为准；LLM-as-Judge 用于快速过滤和趋势监控。
\end{itemize}


\medskip
\hrule
\medskip



\markboth{附录}{附录}


\section{附录：面试答题框架建议}


在面试中回答 AI Agent 相关问题，建议采用以下框架：


\begin{enumerate}[itemsep=1pt, topsep=2pt]
  \item \textbf{先给定义}：一句话说清楚这个技术是什么；
  \item \textbf{再讲原理}：为什么需要它，解决什么问题；
  \item \textbf{对比分析}：和相近技术有什么区别（如 RAG vs Memory | MCP vs Function Calling | Agent vs Workflow）；
  \item \textbf{工程实践}：怎么落地，注意什么坑；
  \item \textbf{Tradeoff}：有什么取舍，什么场景适用，什么场景不适用。
\end{enumerate}


\textbf{一条重要的面试金句}：Agent = Model + Harness。模型负责推理和生成，Harness 负责环境、工具、反馈、沙箱、权限、观测和恢复。决定 Agent 表现上限的，往往不是模型有多强，而是 Harness 有多稳。


\end{document}
